\scanpage{181}
Il en est encore de \unsure{$\overline{\mathrm m}$|même} pour le troisième, si on suppose que $X$ est \unsure{qu.|quasi} cpct séparé sur $S$ et que $\underline H_f^0(\underline O_X)=\underline O_S$ --- de sorte que la $S$-équivalence est identique à la Picard-équivalence rel. à $Y$

Supposons dorénavant que

$X$ qu-compact séparé sur $S$, $\underline H_f^0(\underline O_X)=\underline O_S$ et Kif kif par extension de la base

et que $\underline{\underline{\operatorname{Pic}}}_{X/S}$ et $\underline{\underline{\operatorname{Div}}}_{X/S}$ existent.

Alors l'ensemble précédent
\[
H^0(S,\underline H_f^0(\underline L)^*/\underline O_S^*)
\]
peut s'interpréter comme l'ensemble des sections $\theta$ de $\underline{\underline{\operatorname{Div}}}$ sur $S$ telles que
\diagram{\begin{tikzcd}[ampersand replacement=\&,column sep=small,row sep=small]
\underline{\underline{\operatorname{Pic}}}\arrow[dr,"\varphi"']\&\&\underline{\underline{\operatorname{Div}}}\arrow[ll,"p"']\arrow[dl,"\psi"]\\
\&S\&
\end{tikzcd}}
$p\theta=\lambda$, où $\lambda$ est la section de $\underline{\underline{\operatorname{Pic}}}$ sur $S$ qui correspond à $\underline L$. Donc, posons $\underline{\underline{\operatorname{Div}}}^{(L)}=\varphi^{-1}(\lambda(S))$ [sous-\unsure{préschéma} de $\underline{\underline{\operatorname{Div}}}$) c'est l'ens des sections de $\underline{\underline{\operatorname{Div}}}_{X/S}^{L}$ sur $S$. \deleted{\unsure{Donc}} Ceci est compatible avec extension de la base $S'$, donc le foncteur $\underline{\underline{\operatorname{Div}}}_{X/S}^{L}(S')=H^0(S',\underline H_{f'}^0(\underline L')^*/\underline O_{S'}^*)$ en $S'$ est représentable. \deleted{Ceci doit} [On obtient donc une condition nécessaire intéressante, en termes de diviseurs, \unsure{pour|que} l'existence
% archive sheet number: 180
